% Generated draft: requires mathematical and source review.
\documentclass[11pt]{article}
\usepackage[margin=25mm]{geometry}
\usepackage{fontspec}
\setmainfont{Latin Modern Roman}
\usepackage{amsmath,amssymb,mathtools,unicode-math}
\setmathfont{Latin Modern Math}
\usepackage{enumitem,needspace,xurl,hyperref}
\setlength{\parindent}{0pt}
\setlength{\parskip}{5pt}
\setlength{\emergencystretch}{3em}
\newcommand{\N}{\mathbb{N}}
\newcommand{\Z}{\mathbb{Z}}
\newcommand{\Q}{\mathbb{Q}}
\newcommand{\R}{\mathbb{R}}
\newcommand{\C}{\mathbb{C}}
\newcommand{\F}{\mathbb{F}}
\newcommand{\A}{\mathbb{A}}
\newcommand{\PP}{\mathbb{P}}
\newcommand{\E}{\mathbb{E}}
\newcommand{\eps}{\varepsilon}
\newcommand{\dd}{\,\mathrm{d}}
\begin{document}
\section*{UnsolvedMath ID: 30006613}
\textbf{Source problem code:} catalogue entry
\subsection{Definitions and mathematical statement}
Let \(X\) denote the set of alternatives. The countably infinite case means \(|X|=\aleph_0\). The supplied excerpt does not specify the preference data on \(X\), the definition of single-peakedness being used, or the split system underlying its Buneman graph.

For terminology, in a connected graph \(G\) with graph metric \(d_G\), the interval between vertices \(u,v\) is
\[
I_G(u,v)=\{w:d_G(u,w)+d_G(w,v)=d_G(u,v)\}.
\]
A median graph is a connected graph for which every three vertices \(a,b,c\) have a unique median, namely the unique element of
\[
I_G(a,b)\cap I_G(b,c)\cap I_G(c,a).
\]
A vertex set is convex if it contains the interval between any two of its vertices. A half-space is a nonempty proper vertex set whose complement and itself are convex. These definitions provide terminology; the excerpt does not identify which half-spaces encode the preference data.

\textbf{Primary questions, with unresolved specifications.} What extension of the referenced Buneman-graph model of single-peaked preference data is possible when \(X\) is countably infinite? Is there a coarse variant of the Buneman graph that models the referenced data and permits computation of medians with controlled error? The curated statement additionally requests accommodation of missing or uncertain half-space assignments. This requirement is retained as a supplied qualification, although those assignments and their uncertainty model are not defined in the quoted excerpt.

A precise construction problem cannot yet be specified: the admissible preference profiles, split system, infinite-graph vertex conditions, representation requirements, metric, intended median operation, computational input model, and error bound are absent. In particular, the excerpt does not determine whether the coarse question is restricted to countably infinite alternatives or is a separate question about the preceding data model.

\textbf{Separate question included in the extracted text.} The excerpt then begins ``Problem 3. Coarse median spaces, Arzhantseva'' and asks whether every coarse median space coarsely embeds into a Hilbert space. A coarse embedding of a metric space \((M,d)\) into a Hilbert space \(H\) is a map \(f:M\to H\) for which there exist nondecreasing functions \(\rho_-,\rho_+:[0,\infty)\to[0,\infty)\), with \(\lim_{t\to\infty}\rho_-(t)=\infty\), such that
\[
\rho_-(d(x,y))\leq \|f(x)-f(y)\|_H\leq \rho_+(d(x,y))
\qquad\text{for all }x,y\in M.
\]
Thus the separate question asks whether such \(H,f,\rho_-,\rho_+\) exist for every coarse median space. The excerpt describes a coarse median space informally as a metric space equipped with a ternary operation satisfying median-algebra axioms up to bounded error; it does not supply the complete quantitative axioms or any additional hypotheses. That informal description is insufficient to fix the quantified class precisely. This adjacent question is preserved separately because the extraction does not establish that it belongs to the primary Buneman-graph problem.
\subsection{Short English statement}
The primary problem asks how Buneman models of single-peaked preferences extend to countably infinite alternatives and whether a coarse model can compute approximate medians, with the curated version also requesting missing or uncertain half-space assignments. Essential data and error specifications are absent. The extraction also contains a separate Hilbert-space coarse-embedding question attributed to Arzhantseva. Both require source clarification and human mathematical review.
\subsection{Sources}
\begin{itemize}
\item Local UnsolvedMath registry: \texttt{data/unsolvedmath/problems.json}, permanent ID \(30006613\), title ``Coarse Buneman Models for Countable Preference Data,'' exact source code \texttt{catalogue entry}. The supplied record lists proposed year 2026 and no proposer for the primary problem.
\item Supplied report citation: \emph{Median Geometry and Applications} (2026), Oberwolfach Reports, pp.~487--540, DOI \texttt{10.4171/owr/2026/8}. \url{https://doi.org/10.4171/owr/2026/8}. The extracted text contains an embedded page header identifying page 530.
\item Supplied literature source, labelled ``OWR report'': \url{https://ems.press/content/serial-article-files/53603}.
\item Literature assessment dated \texttt{2026-08-22}: supplied dated curation, not independently verified. The record labels the problem open and reports that the extracted countable-single-peakedness and coarse Buneman questions are unresolved, that the report gives no solution, and that no later resolution was found. These are supplied assessment claims, not a new literature verification.
\item The adjacent coarse-median passage attributes its question to Arzhantseva and the introduction of coarse median spaces to Brian Bowditch. It lists Gromov hyperbolic spaces, CAT(0) cube complexes, mapping class groups, Teichm\"uller spaces of compact surfaces, and hierarchically hyperbolic spaces as examples. It also states that an affirmative Hilbert-embedding answer would generalize a known positive result for median spaces. These are claims in the supplied excerpt; their metric conventions, hypotheses, and supporting references are not supplied or independently verified here.
\item The registry labels its statement \texttt{corrected\_verified} and says only formula delimiters, notation, grammar, and typography were repaired, with mathematical content and scope preserved. This is supplied editorial metadata and does not resolve the extraction and formulation issues described here.
\end{itemize}
\end{document}
